\section{Conclusions}\label{sec:conclusions}
A workflow for kinetic mechanism identification in fed-batch bioprocesses has been presented that learns the kinetics with a hybrid neural ODE before any mechanism is assumed, and then uses the learned rates to start, propose and check mechanistic rate laws. The main conclusions are as follows.
\begin{itemize}
\item A hybrid neural ODE whose rates are non-negative and vanish without substrate learns the signature of each kinetic law from sparse offline data of fed-batch runs, without numerical differentiation; the substrate constraint is required for substrate-limited operation.
\item Learned rates are a reliable source of starting points for refining a library of candidate mechanisms, but not a substitute for the refinement: selection must be decided by fits to the measurements, and an adequacy test is needed to detect when no candidate explains the data.
\item The information content of the experiment, rather than the fitting, limits identification. One run in which the states change monotonically and in concert left the mechanism ambiguous in a third of the cases; three runs with decorrelated biomass and substrate, fitted jointly, identified the generating mechanism in \ConeAccJoint\,\% of the cases.
\item Symbolic regression restricted to physically admissible rational and exponential laws, and judged on the same criterion as the library, produced no false discoveries, recovered non-rational library laws, and discovered a mechanism missing from the library whenever the data required it; adding the discovered law to the library made it available for later runs without disturbing the identification of the existing mechanisms.
\item On four real xylitol fermentations with unknown measurement errors, every run analysed alone selected the same mechanism, with substrate and product inhibition of the conversion and cell death. Joint fits showed that the runs share this mechanism but not its parameter values, so the shared-kinetics assumption on which joint fitting relies must be tested on real data. The symbolic regression consistently pointed to exponential product inhibition with cell death, but without an adequacy test no proposed law was accepted.

\end{itemize}
The workflow runs in minutes per analysis on a workstation and requires only the mass balances, a library of candidate mechanisms and a dictionary of admissible rate-law terms. Coupling it with the model-based design of discriminating experiments is the natural next step.
